\section{Packing tensors into a subspace of small dimension}
\label{sec:packing}

Selecting an averaging length changes the joint distribution of the
variables being averaged.  The deterministic Hilbert-space estimate in
this section will control that change.  In the application, two independent
groups of variables supply the two $L^2$ spaces whose tensor product we
consider.  The estimate is uniform over the testing subspace: a subspace
of moderately growing dimension captures a substantial part of only a
small fraction of the tensors in an almost orthogonal family.

Throughout this section the Hilbert spaces are real, tensor products are
Hilbert tensor products, and logarithms are natural.

\begin{theorem}[Tensor packing]\label{thm:packing}
There is a universal constant $C$ with the following property.  Let
$0<\varepsilon\le1$, $d\ge1$, and $L\ge0$ be given, with $d,L$ integers.
Let $x_1,\ldots,x_m$ and $y_1,\ldots,y_m$ lie in the unit balls of real
Hilbert spaces $H_1,H_2$.  For each family suppose that a graph on
$\{1,\ldots,m\}$ of maximum degree at most $L$ contains every off-diagonal
pair whose absolute inner product exceeds $\eta$, where
\[
                         0\le\eta\le\varepsilon^2/10^4.
\]
The two exceptional graphs may differ.  For every orthogonal projection
$P$ on $H_1\otimes H_2$ of rank at most $d$,
\begin{equation}\label{eq:packing-quantitative}
 \#\{n:\norm{P(x_n\otimes y_n)}\ge\varepsilon\}
 \le \frac{C(1+L)}{\varepsilon^2}
       \exp\!\left(C\varepsilon^{-2}\sqrt d\log(2+d)\right).
\end{equation}
In particular, for
\begin{equation}\label{eq:packing-dimension}
                  d_m=\left\lceil(1+\log m)^{3/2}\right\rceil,
\end{equation}
the left side is $o_{\varepsilon,L}(m)$, uniformly over the Hilbert
spaces, the vectors, the projections, and the allowed values of $\eta$.
\end{theorem}

The proof separates directions on which the testing space has large
variance, then applies a Gaussian comparison to the remaining directions.
The comparison is the finite form of the Sudakov--Fernique inequality;
see \cite[Section~1, Equation~(1.2)]{Fernique} and the minoration in
\cite{Sudakov}.  We first record the matrix estimate needed for this
application.  Its proof is the Gaussian integration-by-parts moment
argument underlying matrix Khintchine estimates; compare
\cite[Lemma~6.1 and Proposition~7.3]{Tropp}.  We include the derivation
to display the dependence on the covariance trace rather than on the
ambient matrix dimensions.

\begin{lemma}[A Gaussian matrix estimate]\label{lem:packing-gaussian}
Let $B_1,\ldots,B_r$ be matrices between two finite-dimensional real
Hilbert spaces, and put
\[
 H_j=\begin{pmatrix}0&B_j\\ B_j^*&0\end{pmatrix},
 \qquad S=\sum_{j=1}^r H_j^2.
\]
If $\norm S_{\rm op}\le\sqrt d$ and $\tr S\le2d$ for some $d\ge1$, then
for independent standard normal variables $g_j$,
\begin{equation}\label{eq:packing-gaussian-bound}
       \E\left\|\sum_{j=1}^r g_jB_j\right\|_{\rm op}
             \le C d^{1/4}\sqrt{\log(2+d)}.
\end{equation}
The constant does not depend on the dimensions of the two spaces.
\end{lemma}

\begin{proof}
Write $H=\sum_jg_jH_j$.  Gaussian integration by parts, followed by
differentiation of the matrix power, gives for $q\ge2$
\[
 \E\tr H^{2q}
  =\sum_{j=1}^r\sum_{a=0}^{2q-2}
       \E\tr\bigl(H_jH^aH_jH^{2q-2-a}\bigr).
\]
For a fixed self-adjoint matrix $H$ with eigenvalues $\lambda_u$, and a
self-adjoint matrix $K$, the absolute value of a summand is at most
\[
 \sum_{u,v}|K_{uv}|^2\,
             |\lambda_v|^a|\lambda_u|^{2q-2-a}
       \le \tr\bigl(K^2|H|^{2q-2}\bigr).
\]
The last inequality is weighted arithmetic--geometric mean, together
with the symmetry $|K_{uv}|=|K_{vu}|$.  Summing in $j$ therefore yields
\begin{equation}\label{eq:packing-moment-recursion}
 \E\tr H^{2q}
       \le(2q-1)\norm S_{\rm op}\,
                         \E\tr |H|^{2q-2}.
\end{equation}
Since $\E\tr H^2=\tr S$, iteration gives
\[
 \E\tr H^{2q}
       \le(2q-1)!!\,\norm S_{\rm op}^{q-1}\tr S.
\]
The operator norm of $H$ equals that of $\sum_jg_jB_j$.  Taking a
$2q$-th root and using the hypotheses bounds its expectation by
\[
             C\sqrt q\,d^{1/4}(2\sqrt d)^{1/(2q)}.
\]
Choose an integer $q\ge2$ comparable to $\log(2+d)$.  This proves
\eqref{eq:packing-gaussian-bound}.
\end{proof}

\begin{proof}[Proof of Theorem~\ref{thm:packing}]
We use repeatedly the following elementary observation.  If unit-ball
vectors $v_n$ have inner products bounded in absolute value by $\eta$
outside a graph of degree $L$, then for any nonempty set $A$ of indices,
\begin{equation}\label{eq:packing-cluster}
 \left\|\frac1{|A|}\sum_{n\in A}v_n\right\|^2
                    \le\frac{1+L}{|A|}+\eta.
\end{equation}
Indeed, there are at most $(1+L)|A|$ ordered diagonal or exceptional
pairs in the expansion of the squared norm.  Thus a lower bound on the
norm of such an average bounds the number of vectors in the set.

It suffices to work in the finite spans of the $x_n$ and $y_n$.
To justify this reduction, let $Q$ project onto their tensor product
and replace $\operatorname{Ran}P$ by $Q(\operatorname{Ran}P)$.
For a vector $z$ in that tensor product, its norm of projection onto
the new subspace is at least $\norm{Pz}$: in the supremum defining
$\norm{Pz}$, a unit testing vector can be replaced by its $Q$-image,
whose norm is at most one.  The dimension does not increase.

Identify tensors with matrices from $H_2$ to $H_1$, and let
$A_1,\ldots,A_r$ be a Hilbert--Schmidt orthonormal basis of the testing
space, where $r\le d$.  The matrices
\[
                  R_1=\sum_jA_jA_j^*,\qquad
                  R_2=\sum_jA_j^*A_j
\]
have trace $r$.  Let $E_1,E_2$ be their spectral subspaces for eigenvalues
strictly greater than $\sqrt d$.  Each has dimension at most $\sqrt d$.

Discard the indices for which the projection of $x_n$ onto $E_1$ has
norm greater than $\varepsilon/10$, and do the same for $y_n,E_2$.
Cover either exceptional subspace's unit ball by at most
$(1+80/\varepsilon)^{\sqrt d}$ balls of radius $\varepsilon/40$.
If a ball contains $s$ of the projections being discarded, their
average has norm at least $\varepsilon/20$.  Applying
\eqref{eq:packing-cluster} to the original vectors, with
$\eta\le\varepsilon^2/10^4$, gives
$s\le C(1+L)\varepsilon^{-2}$.  The number discarded is consequently
at most
\begin{equation}\label{eq:packing-discarded}
 C(1+L)\varepsilon^{-2}
              \exp\bigl(\sqrt d\log(1+80/\varepsilon)\bigr).
\end{equation}

Let $Q_1,Q_2$ project onto $E_1^\perp,E_2^\perp$, and set
$B_j=Q_1A_jQ_2$.  The coefficient map
\[
 z\longmapsto (\ip{A_j}{z}_{\rm HS})_{j=1}^r
\]
is a contraction.  For an index that remains,
\[
 \norm{x_n\otimes y_n-Q_1x_n\otimes Q_2y_n}
   \le\norm{(1-Q_1)x_n}+\norm{(1-Q_2)y_n}
   \le\varepsilon/5.
\]
Hence every remaining index counted on the left of
\eqref{eq:packing-quantitative} has coefficient vector
\[
                  v_n=(\ip{x_n}{B_jy_n})_{j=1}^r,
                  \qquad \norm{v_n}\ge\varepsilon/2.
\]
The map from $x_n\otimes y_n$ to $v_n$ is still contractive.
Outside the union of the two exceptional graphs, the tensor inner
products have absolute value at most $\eta^2$; the union has degree at
most $2L$.  A set of these coefficient vectors lying in a ball of
radius $\varepsilon/8$ centred at one of them has average norm at
least $3\varepsilon/8$.  Applying \eqref{eq:packing-cluster} to their
preimage tensors shows that the set has at most
$C(1+L)\varepsilon^{-2}$ indices.  A maximal
$\varepsilon/8$-separated subset therefore has cardinality $l$ such
that the number still to be counted is at most
$C(1+L)\varepsilon^{-2}l$.

It remains to bound this separated set.  For
$G=\sum_jg_jB_j$, the centered Gaussian variables
$Z_n=\ip{x_n}{Gy_n}$ have increment variances
$\E|Z_n-Z_{n'}|^2=\norm{v_n-v_{n'}}^2$.
The Gaussian comparison just mentioned gives, for $l\ge2$,
\begin{equation}\label{eq:packing-comparison}
 c\varepsilon\sqrt{\log l}
                     \le\E\max_nZ_n\le\E\norm G_{\rm op}.
\end{equation}
For completeness, the comparison of two finite centered Gaussian
families follows by applying Gaussian integration by parts to
$\beta^{-1}\log\sum_n e^{\beta z_n}$ along their independent Gaussian
interpolation.  Its derivative is
\[
 \frac\beta4\E\sum_{n,n'}p_np_{n'}
   \left(\E|Z_n-Z_{n'}|^2-\E|Z'_n-Z'_{n'}|^2\right),
 \qquad p_n=\frac{e^{\beta z_n}}{\sum_j e^{\beta z_j}},
\]
and is nonnegative when the first increment variances dominate.
Letting $\beta$ tend to infinity compares the expected maxima.
In \eqref{eq:packing-comparison}, take $Z'_n$ to be independent normals
of variance $\varepsilon^2/128$ and use the elementary lower bound
$\E\max_{n\le l}g_n\ge c\sqrt{\log l}$.

Compression gives
\[
 \sum_jB_jB_j^*\le Q_1R_1Q_1,\qquad
 \sum_jB_j^*B_j\le Q_2R_2Q_2.
\]
Their norms are at most $\sqrt d$, and their traces are at most $d$.
Lemma~\ref{lem:packing-gaussian} applies, so
\eqref{eq:packing-comparison} implies
\[
                   \log l\le C\varepsilon^{-2}\sqrt d\log(2+d).
\]
The case $l\le1$ satisfies the resulting bound as well.  Combining
this with \eqref{eq:packing-discarded}, and absorbing
$\log(1+80/\varepsilon)\le C\varepsilon^{-2}$, proves
\eqref{eq:packing-quantitative}.

For $d=d_m$, the exponent in that estimate is
$O_\varepsilon((\log m)^{3/4}\log\log m)=o(\log m)$.
The asserted $o(m)$ conclusion follows.
\end{proof}

We will use the conclusion also in averaged form.  Since every
projection norm is at most one, for any fixed accuracy $\varepsilon$,
vectors satisfying the corresponding small-pair hypothesis, and an
orthogonal projection $P$ of rank at most $d_m$,
\begin{equation}\label{eq:packing-average}
 \frac1m\sum_{n=1}^m\norm{P(x_n\otimes y_n)}
   \le\varepsilon+
       \frac{C(1+L)}{m\varepsilon^2}
       \exp\!\left(C\varepsilon^{-2}\sqrt{d_m}\log(2+d_m)\right).
\end{equation}
Here $\varepsilon$ and then $\eta$ are fixed before $m$ increases.
The exceptional degree may depend on this chosen $\eta$.  In
Section~\ref{sec:channels}, mixing supplies precisely this order of
quantifiers; a uniform rate of mixing will not be needed.
